\section{Wstęp}

\subsection{Geneza}

Pomysł na realizację niniejszego projektu zrodził się podczas analizowania wysokopoziomowych testów modułowych. Wspomniane testy charakteryzują się wysokim poziomem abstrakcji oraz szerokim zakresem funkcjonalności, co skutkuje dużą liczbą wymaganych powtarzalnych ustawień. Na tej podstawie zaobserwowano interesującą tendencję polegającą na tym, że znaczna część kodu testowego okazywała się redundantna, a wiele mechanizmów było wielokrotnie powielanych poprzez kopiowanie oraz wprowadzanie drobnych modyfikacji do istniejących fragmentów kodu.

Naturalnym rozwiązaniem tego problemu mogłoby wydawać się wydzielanie powtarzalnych elementów do funkcji pomocniczych lub bibliotek testowych oraz utrzymywanie ich w uporządkowanej i współdzielonej formie. Podejście to jednak wiąże się z istotnymi trudnościami, zarówno natury organizacyjnej, jak i technicznej. W praktyce często prowadzi ono do nadmiernego rozrostu bibliotek pomocniczych, niejednoznacznych interfejsów oraz problemów z ich długoterminowym utrzymaniem.

W szczególności pojawiają się następujące problemy:
\begin{enumerate}
  \item Jak określić moment, w którym dana funkcjonalność jest wystarczająco generyczna, aby mogła zostać wydzielona do wspólnej biblioteki?
  \item Jak zminimalizować ryzyko tworzenia funkcji lokalnie w plikach testowych, które z czasem zaczynają być wykorzystywane przez inne testy, prowadząc do niejawnych zależności pomiędzy teoretycznie niezależnymi modułami?
  \item Jak zapewnić spójny sposób tworzenia oraz utrzymania takich funkcji w projektach rozwijanych przez dziesiątki, a niekiedy setki programistów, o różnym poziomie doświadczenia i odmiennych nawykach pracy?
\end{enumerate}

Jednocześnie, jak powszechnie wiadomo, modele sztucznej inteligencji uczą się na podstawie zbiorów danych treningowych, a następnie wykorzystują wyuczone wzorce do rozwiązywania nowych problemów. Obserwacja ta stała się punktem wyjścia do rozważenia możliwości zastosowania modeli sztucznej inteligencji (ang. Artificial Intelligence, AI) w obszarze automatycznego generowania wspólnych, powtarzalnych fragmentów kodu testowego na podstawie już istniejących przykładów.

Jednakże, aby taki system mógł zostać efektywnie wdrożony w praktyce, niezbędne jest zaprojektowanie i utrzymanie procesu ciągłej integracji i ciągłego dostarczania (ang. Continuous Integration / Continuous Delivery, CI/CD), zapewniającego automatyczne trenowanie modeli, ich ewaluację oraz kontrolowane wdrażanie w środowisku produkcyjnym. W przeciwnym razie konieczne byłoby wyznaczenie osoby odpowiedzialnej za manualne utrzymanie i nadzór.

W ten sposób wyłaniają się dwa kluczowe pytania stanowiące oś niniejszego projektu:
\begin{enumerate}
\item Czy stworzenie modelu AI umożliwiającego automatyczne generowanie fragmentów kodu testowego jest rozwiązaniem możliwym do realizacji oraz uzasadnionym ekonomicznie?
\item Czy bardziej opłacalne jest oddelegowanie pracownika do ręcznego utrzymywania i rozwijania tego rozwiązania, czy też zasadne jest jego zautomatyzowanie z wykorzystaniem narzędzi DevOps oraz procesów CI/CD?
\end{enumerate}

\subsection{Założenia}

Ze względu na złożoność problemu generowania kodu testowego oraz ograniczony czas realizacji projektu zdecydowano się na zastosowanie problemu zastępczego, umożliwiającego weryfikację postawionej hipotezy bez konieczności budowania zaawansowanego modelu językowego. Jako problem zastępczy wybrano segmentację semantyczną obrazów wodomierzy z wykorzystaniem architektury U-Net. Wybór ten jest uzasadniony, ponieważ modele segmentacyjne posiadają jednoznacznie zdefiniowane metryki jakości, takie jak współczynnik Dice (ang. Dice coefficient, miara nakładania masek) oraz IoU (ang. Intersection over Union, stosunek części wspólnej do sumy obszarów), charakteryzują się umiarkowanymi wymaganiami obliczeniowymi umożliwiającymi trening w ramach dostępnego budżetu, a także opierają się na dostępnych modelach wstępnie wytrenowanych, co pozwala na obiektywne porównanie kolejnych wersji.

Dla uproszczenia przyjęto założenie, że model sztucznej inteligencji dostarczany jest przez zewnętrzną firmę, natomiast zakres pracy obejmuje jego dalsze uczenie, ocenę jakości oraz wdrażanie. Takie podejście odzwierciedla realny scenariusz środowiska produkcyjnego, w którym zespół DevOps otrzymuje gotowy model i odpowiada za zarządzanie jego cyklem życia.

W tak zdefiniowanym kontekście drugie pytanie sformułowane w poprzedniej sekcji pozostaje aktualne: czy bardziej opłacalne jest oddelegowanie pracownika do ręcznego utrzymywania i rozwijania tego rozwiązania, czy też zasadne jest jego zautomatyzowanie z wykorzystaniem narzędzi DevOps oraz procesów CI/CD? W podejściu manualnym proces ponownego uczenia i wdrażania modelu obejmuje ręczne przygotowanie danych, uruchamianie treningu, walidację wyników oraz publikację nowej wersji. Podejście to, mimo swojej prostoty koncepcyjnej, wiąże się z ryzykiem błędów ludzkich, trudnościami w zachowaniu powtarzalności eksperymentów oraz ograniczoną skalowalnością. Automatyzacja natomiast umożliwia jednoznaczne powiązanie wersji modelu z wersją danych treningowych, wprowadzenie obiektywnej bramki jakości oraz eliminację ręcznych etapów mogących prowadzić do niespójności.

W ramach niniejszej pracy opracowano kompletny, zautomatyzowany system CI/CD dla modeli sztucznej inteligencji. System obejmuje wersjonowanie danych treningowych z wykorzystaniem narzędzia DVC, śledzenie eksperymentów i rejestr modeli (MLflow), automatyczny potok (and. pipeline) treningowy uruchamiany w GitHub Actions, mechanizm bramki jakości porównujący nowy model z aktualną wersją produkcyjną oraz wdrożenie modelu w środowisku chmurowym za pomocą lekkiej dystrybucji Kubernetes (k3s), uruchomionej na instancji Amazon EC2  (ang. Elastic Compute Cloud, Elastyczna Chmura Obliczeniowa) wraz z mechanizmami monitoringu (Prometheus, Grafana). Całość zaprojektowano w taki sposób, aby dodanie nowych danych treningowych wymagało jedynie wykonania polecenia \texttt{git push}, a pozostałe etapy procesu przebiegały automatycznie. Tak zdefiniowany zakres umożliwia bezpośrednie porównanie kosztów i nakładu pracy podejścia manualnego oraz zautomatyzowanego, co stanowi główny cel analityczny pracy.
